\documentclass[12pt]{article}
\usepackage[utf8]{inputenc}
\usepackage[T1]{fontenc}
\usepackage{hyperref}
\title{Toku-No-Shima Language Resources for K-12 Curriculum Development}
\author{Compiled by Claude Sonnet 5 (claude-sonnet-5) for Andr\'as Kornai}
\date{}
\begin{document}
\maketitle

\section{Introduction}
Toku-No-Shima (Tokunoshima, ISO 639-3 tkn) is a Japonic (Amami--Okinawan) language spoken on Tokunoshima, Kagoshima Prefecture, Japan. Speakers call it Shimaguchi. Wikipedia gives about 5,100 native speakers (2004). Search found little material designed for schools. What exists is mostly academic.

\section{Primary Educational Resources}
No primers, textbooks or school programmes were found in the search results.

\section{Linguistic and Reference Materials}
\begin{itemize}
\item Tokunoshima language, Wikipedia overview: \url{https://en.wikipedia.org/wiki/Tokunoshima_language}
\item ``A Study for Compilation of Tokunoshima Dialect Dictionary'' (KAKEN research project record): \url{https://kaken.nii.ac.jp/grant/KAKENHI-PROJECT-15320050/}
\item Paper on building verb and adjective paradigms from a dialect corpus (the Sawaki group's dictionary work, based on a Tokunoshima translation of Frei's 2000 sentences): \url{https://journals.um.si/index.php/slaviacentralis/en/article/view/653}
\item Sawaki, Nakajima and Fukushima, NWAVAP 2 presentation (2013) on verb forms and the loss of honorific forms among younger speakers: \url{https://www2.ninjal.ac.jp/past-events/nwavap02/SawakiNakajimaFukushima-NWAVAP2-2013.pdf}
\item Hosei University index listing Uchima's \emph{Grammar of the Tokunoshima Inokawa Dialect} and Nakamoto's Inokawa vocabulary (1979): \url{https://www.en.hosei.ac.jp/LUC2HOSEI/cdata/luc2hosei_19504_jaen.html}
\item Repository record (Japanese academic paper, content not verified): \url{https://unii.repo.nii.ac.jp/records/216}
\end{itemize}

\section{Religious and Community Resources}
None found.

\section{Digital and Audio Resources}
The dictionary project above is described as multimedia. No public audio or online dictionary URL was found in the results.

\section{Curriculum Development Context}
The language is spoken as a first language by adults only and is not known to be taught in schools. Any curriculum would have to be built from scratch, drawing on the academic corpus and dictionary work above. The next step is to contact the researchers behind that work and look for Japanese-language local materials, which this search did not cover.
\end{document}